\documentclass[../../main.tex]{subfiles}
\begin{document}
\section{Introduction and main conditional statements}
\label{sec:introduction}

Let $\mu$ be an isotropic log-concave probability measure on $\R^n$:
\[
  \E_\mu X=0,\qquad \Cov_\mu(X)=I_n .
\]
Write $\mu^+(E)$ for the Minkowski boundary measure of a Borel set $E$, and let
\[
  I_\mu(p)=\inf\{\mu^+(E):\mu(E)=p\},\qquad 0<p<1,
\]
be the isoperimetric profile. The Cheeger constant is
\[
  h_\mu=\inf_{0<p<1}\frac{I_\mu(p)}{\min(p,1-p)}.
\]
The KLS conjecture asserts that $h_\mu\ge c>0$ for every isotropic log-concave $\mu$, with $c$ a
universal constant \cite{KannanLovaszSimonovits1995,LeeVempala2018}. Equivalently, up to universal constants in
this category, the Poincare constant should be bounded independently of $n$. We denote the inverse
Cheeger scale by
\[
  \PsiKLS_\mu=h_\mu^{-1}.
\]
This avoids the common ambiguity in the literature, where the symbol $\psi_\mu$ may denote either a
Cheeger constant or its reciprocal. We also write
\begin{equation}\label{eq:hstar-def}
  \hstar_n=\inf\bigl\{h_\nu:\ \nu\ \text{isotropic log-concave on }\R^n\bigr\},
\end{equation}
the worst-case Cheeger constant of the dimension; the KLS conjecture is the assertion
$\inf_n\hstar_n>0$. The published benchmark is $\hstar_n\ge c(\log n)^{-1/2}$
\cite{Klartag2023Logarithmic}. A July 2026 version-1 preprint of Letwin improves the claimed
bound to $\hstar_n\ge c(\log n)^{-1/4}$ \cite{Letwin2026QuadraticKLS}; throughout this report,
uses of that improvement are explicitly conditional on the preprint input.

The surrounding landscape has changed substantially. Letwin's moment-map/Stein-kernel argument
proves a dimension-free quadratic Poincar\'e inequality for quadratic forms and then combines it
with the log-concave Lichnerowicz comparison to obtain the preceding $\log^{1/4}n$ inverse-Cheeger
bound \cite{Letwin2026QuadraticKLS}. Eldan's stochastic localization reduced
thin-shell information to spectral-gap information up to logarithmic factors \cite{Eldan2013ThinShell}, and
Lee--Vempala developed this into powerful isoperimetric, concentration, and mixing estimates
\cite{LeeVempala2024}, in a tradition that also includes Paouris-type concentration bounds
\cite{Paouris2006}. Klartag proved that KLS and slicing hold up to a factor $\sqrt{\log n}$ by an
improved log-concave Lichnerowicz inequality \cite{Klartag2023Logarithmic}. Bourgain's slicing problem has
since been resolved through the Klartag--Lehec argument using Guan's covariance-process bound,
Milman's $M$-ellipsoid theory, and Shannon--Stam stability
\cite{Guan2024,MilmanSchechtman1986,EldanMikulincer2020,KlartagLehec2025Slicing}. The thin-shell
conjecture has also been resolved by Klartag--Lehec through parallel coupling and stochastic
localization \cite{KlartagLehec2025ThinShell}. These results remove major global radial and volume
obstructions. They do not, by themselves, control a fixed bottleneck set under localization.

The guiding principle of this report is therefore set-specific:
\[
  \boxed{\text{follow the posterior evolution of a fixed candidate cut }E.}
\]
The mass $p_t=\mu_t(E)$ is a martingale. KLS follows if a balanced near-bottleneck cut cannot be
identified too rapidly by the early Gaussian observation generated by localization.

\subsection{The two proof routes}

The report separates two routes.

\begin{enumerate}[label=\textbf{Route \Alph*.},leftmargin=2.5em]
  \item \textbf{All-cut stochastic route.}  Prove an absorptive two-color Carleson estimate for
  every balanced cut. This immediately yields KLS.

  \item \textbf{Near-Cheeger geometric route (weighted form).}  Work only with near-minimizers of
  the isoperimetric profile. This is sufficient for KLS, but it requires two additional
  mechanisms, which must be taken in \emph{covariance-weighted}
  form: weighted excess propagation under localization, and a weighted stable Stein-trace
  estimate. The consumption audit of Section~\ref{sec:excess} shows that the unweighted form of
  the second mechanism implied KLS by itself, with the unweighted excess term
  inert; the quantified two-tail obstruction shows that the unweighted form admits no slice-wise
  proof. The weighted form stated below is the correct formulation of the route.
\end{enumerate}

The second route is the natural one for the boundary/Jacobi ideas. It should not be advertised as
an all-cut estimate: arbitrary balanced sets can have enormous two-color covariance contrast even
when their perimeter is small in an anisotropic posterior.

\subsection{Main stochastic quantities}

For orientation we recall the two-color quantities; they are defined in full in
Section~\ref{sec:notation}. Fix a measurable set $E$ and put $F=E^c$. Under localization define
\[
  p_t=\mu_t(E),\qquad q_t=1-p_t,
  \qquad s_t=p_tq_t,
\]
\[
  \delta_t=m_t^E-m_t^F,
  \qquad G_t=\Sigma_t^E-\Sigma_t^F,
  \qquad r_t=s_t\abs{\delta_t}^2,
\]
where $m_t^E,m_t^F$ and $\Sigma_t^E,\Sigma_t^F$ are the two conditional means and covariances. The
Riccati source and damping are
\[
  S_t=s_t\norm{G_t}_{\HS}^2,
  \qquad
  D_t=2s_t\delta_t^TA_t\delta_t-r_t^2,
\]
where $A_t=\Cov(\mu_t)$. The isoperimetric excess process is
\[
  e_t(E)=\mu_t^+(E)-I_{\mu_t}(p_t)\ge0,
\]
and the covariance-excess functionals, central to both routes, are
\begin{equation}\label{eq:interface-def}
  X_t=\bigl(\lmax(A_t)-1\bigr)_+,
  \qquad
  \Xi_T(\mu)=\int_0^T\E X_t\dd t .
\end{equation}
Let $\tau$ be the first time at which $p_t$ exits a fixed balanced interval, for instance
$[1/3,2/3]$, and let $\tau_\eta$ be the tight window of \eqref{eq:tau-tight} below.

\subsection{Main conditional statements}

\begin{assumption}[All-cut absorptive two-color Carleson estimate]
\label{ass:all-cut-carleson}
There exist universal constants $T_0,C_0,C_1$ and $\alpha<1$ such that for every isotropic
log-concave $\mu$, every balanced measurable set $E$, and every interval $I\subset[0,T_0]$,
\begin{equation}\label{eq:all-cut-carleson}
  \E\int_{I\cap[0,\tau]} S_t\dd t
  \le
  C_0\abs I+C_1\E\int_{I\cap[0,\tau]}r_t\dd t
  +\alpha\E\int_{I\cap[0,\tau]}D_t\dd t .
\end{equation}
\end{assumption}

\begin{theorem}[All-cut Carleson implies KLS]
\label{thm:intro-all-cut}
Assumption~\ref{ass:all-cut-carleson} implies the KLS conjecture.
\end{theorem}

The proof is given in Sections~\ref{sec:mass-martingale}--\ref{sec:carleson}. It uses only
stochastic localization, the two-color Riccati identity, and the fact that a $T$-uniformly
log-concave posterior has a dimension-free Cheeger lower bound at scale $\sqrt T$.

For the geometric route, the central package is now stated in weighted form.

\begin{assumption}[Weighted near-Cheeger geometric package]
\label{ass:weighted-package}
There are universal constants $T_0,C_0,C_1,C_2$, $0\le\beta<\tfrac12$, $\gamma>0$, and
$\eta\in(0,\tfrac14]$ satisfying
\begin{equation}\label{eq:weighted-absorption-margin}
  2\beta+64\eta^2<1,
\end{equation}
such that, for the corresponding fixed tight stopping window $\tau_\eta$, the following two
estimates hold for every initial
near-Cheeger set $E$ with $p_0=1/2$ (or $\abs{p_0-1/2}\le\eta/2$).

\begin{enumerate}[label=(\roman*-w),leftmargin=3.2em]
  \item \textbf{Weighted excess propagation.}
  \begin{equation}\label{eq:intro-weighted-excess}
    \E\int_0^{T\wedge\tau_\eta} e_t(E)\,\bigl(1+\norm{A_t}_\op\bigr)^{5/2}\dd t
    \le C_2\bigl(T e_0(E)+T^{1+\gamma}\bigr),
    \qquad 0<T\le T_0 .
  \end{equation}

	  \item \textbf{Weighted stable Stein-trace estimate.}  With the Stein-trace norm
	  $\calS_{\mu_t}(E)$ defined in Section~\ref{sec:stein},
	  \begin{equation}\label{eq:intro-weighted-stein}
	  \begin{aligned}
	    \E\int_0^{T\wedge\tau_\eta}\frac{\calS_{\mu_t}(E)}{s_t}\dd t
	    &\le
	    C_0T+C_1\E\int_0^{T\wedge\tau_\eta}r_t\dd t
	    +\beta\E\int_0^{T\wedge\tau_\eta}D_t\dd t \\
	    &\qquad
	    +C_2\E\int_0^{T\wedge\tau_\eta}
	    e_t(E)\bigl(1+\norm{A_t}_\op\bigr)^{5/2}\dd t .
	  \end{aligned}
	  \end{equation}
\end{enumerate}
\end{assumption}

\begin{theorem}[Weighted geometric package implies KLS]
\label{thm:intro-weighted}
Assumption~\ref{ass:weighted-package} implies the KLS conjecture.
\end{theorem}

The proof is given in Section~\ref{sec:stein}, using the audit and propagation analysis of
Section~\ref{sec:excess}. The weight $(1+\norm{A_t}_\op)^{5/2}$ is calibrated on the anisotropic
two-tail configuration (Proposition~\ref{prop:two-tail}): among pure powers of $\lmax(A_t)$
multiplying absolute excess in this slice-wise package, $5/2$ is the smallest statically
consistent exponent. This does not assert that every possible proof of KLS needs this weight.
The following audit explains the constraint on the present formulation.

\begin{proposition}[Consumption audit; deceptive strength of the unweighted estimate]
\label{prop:intro-audit}
\leavevmode
\begin{enumerate}[label=(\alph*),leftmargin=2.2em]
\item For every isotropic log-concave $\mu$, every measurable $E$ with $\mu(E)=1/2$ and excess
$e_0=e_0(E)$, every $T>0$ and every stopping time $\tau$,
\[
  \E\int_0^{T\wedge\tau}e_t(E)\dd t\ \le\ (1+e_0)\,T .
\]
\item Consequently, the \emph{unweighted} stable Stein-trace estimate --- estimate
\eqref{eq:intro-weighted-stein} with the weight $(1+\norm{A_t}_\op)^{5/2}$ replaced by $1$ ---
implies the KLS conjecture by itself; the unweighted excess term is inert, and unweighted excess
propagation is unconditionally true in the form consumed.
\item Nevertheless, the unweighted estimate admits no slice-wise proof: there are log-concave
pairs $(\nu,E)$ with $\nu(E)=1/2$, $r=D=0$, excess $e(E)\to0$, and $\calS_\nu(E)/s\to\infty$
(Proposition~\ref{prop:two-tail}).
\end{enumerate}
\end{proposition}

The proof is given in Section~\ref{sec:excess}. Item (b) is not good news about excess
propagation; it is a warning that the unweighted geometric package concentrated its entire
logical weight, invisibly, in the Stein-trace estimate. Item (c) shows that this weight cannot be
discharged one time-slice at a time. The weighted formulation of
Assumption~\ref{ass:weighted-package} redistributes the weight correctly: the weighted excess term
is no longer inert, and the static obstruction is consistent with it.

Finally, Section~\ref{sec:bootstrap} provides the propagation mechanism available today: a
bootstrap comparison theorem anchored at $\hstar_n$, valid for \emph{every} balanced cut of a
near-worst measure --- the only regime a proof of KLS by contradiction requires --- which
compresses excess propagation into the scalar interface functional $h_\mu\,\Xi_T(\mu)$, together
with a complete evaluation of that interface against the known covariance technology.

\subsection{Status of the ingredients}

\begin{center}
\begin{tabularx}{\textwidth}{@{}>{\raggedright\arraybackslash}p{0.30\textwidth}>{\raggedright\arraybackslash}p{0.31\textwidth}Y@{}}
\toprule
Ingredient & Role & Status in this report \\
\midrule
Mass martingale and stopped centroid reduction & Reduces KLS to survival of a balanced cut for universal time & Proved. \\
\addlinespace
Matrix and scalar two-color Riccati identities & Isolate the source $S_t$ and damping $D_t$ exactly & Proved formally, with standard approximation conventions. \\
\addlinespace
Per-direction Carleson estimate & Shows that $s_tG_t^2$ has dimension-free occupation in every fixed direction & Proved from the matrix Riccati identity. \\
\addlinespace
Operator-to-trace upgrade & The true stochastic missing estimate & Open. \\
\addlinespace
Quadratic-chaos thin shell & Static time-zero form of two-color covariance control & Imported with constant $8$ from a July 2026 version-1 preprint; whitening still leaves dynamic covariance alignment. \\
\addlinespace
Consumption audit of the excess term & Shows the unweighted excess term is inert and the unweighted Stein-trace estimate alone implies KLS & Proved (Section~\ref{sec:excess}). \\
\addlinespace
Quantified two-tail obstruction and weight calibration & Forces the covariance weight $(1+\norm{A_t}_\op)^{5/2}$ on any slice-wise excess term; the $T^{1+\gamma}$ rate remains a strong open target & Proved (Section~\ref{sec:qcts}). \\
\bottomrule
\end{tabularx}
\end{center}

\begin{center}
\begin{tabularx}{\textwidth}{@{}>{\raggedright\arraybackslash}p{0.30\textwidth}>{\raggedright\arraybackslash}p{0.31\textwidth}Y@{}}
\toprule
Ingredient & Role & Status in this report (continued) \\
\midrule
Perimeter martingale and exact excess identity & Reduces propagation to the localized profile floor; exposes a circularity risk without proving a no-go theorem & Proved (Section~\ref{sec:excess}). \\
\addlinespace
Bootstrap comparison theorem & Reduces near-worst excess propagation to the interface $h_\mu\,\Xi_T$ & Proved (Section~\ref{sec:bootstrap}). \\
\addlinespace
Interface evaluation & $\log\log n$ loss from known covariance technology; crude evaluation provably insufficient; a relative bound for all measures at a sufficiently small universal time is KLS-sufficient & Proved, given the imported covariance input of Hypothesis~\ref{hyp:KI}. \\
\addlinespace
Weighted excess propagation & Keeps a near-Cheeger cut near-minimal at the weighted scale & Open and central. \\
\addlinespace
Extremality tames the covariance process & Residual target of the bootstrap route & Open; proposed coupling of the splitting analysis to the covariance SDE. \\
\addlinespace
Weighted stable Stein trace & Boundary route to the two-color source estimate & Open; Reilly/Jacobi gives a proposed mechanism, but a uniform almost-stability trace lemma is still missing. \\
\addlinespace
Jacobi constant-mode splitting & Proposed treatment of modes invisible to volume-preserving stability & One sufficient split model is proved; converse rigidity, coercivity modulo zero modes, and quantitative stability are open. \\
\bottomrule
\end{tabularx}
\end{center}
\end{document}
